\section{Bernadette After the Apparitions: Nevers, Death,
Canonization}

\subsection{The hospice years, 1860--1866}

After 1858 Bernadette's position in Lourdes became impossible in a
way modernity recognizes at once: she was a celebrity, and the
crowds came for her. The received exchanges preserve how she
carried it---no payment (``I want to remain poor''), no blessing
of rosaries (``I don't wear a stole''), no selling of medals
(``I'm not a merchant''), and, faced with pictures of herself at
ten sous, ``Ten sous, that's all I'm worth!'' To shelter her, the
parish priest and the mayor agreed to place her with the Sisters
of Nevers at the town hospice, where she was admitted as ``a sick
poor person'' on 15 July 1860; there she finally learned to read
and write, cared for the sick, and weighed a vocation she thought
closed to a poor girl without dowry or health. She chose the
Sisters of Charity of Nevers, in her own reported words,
``because they did not try to attract me.''\endnote{Sanctuary of
Lourdes, ``Bernadette Soubirous''
(\url{https://www.lourdes-france.com/en/bernadette-soubirous/}),
read 2026-07-25; all quoted phrases are the Sanctuary's received
account. The canonization \term{Litterae decretales} add the
canonical sequence: her request to enter was handled from 1864
with the counsel and authority of the Bishop of Nevers, and poor
health held her at the Lourdes hospice two further years
(\work{Acta Apostolicae Sedis} 26 (1934), p.~76).}

\subsection{The hidden life at Nevers, 1866--1879}

In July 1866 she left Lourdes---taking leave of the grotto she
would never see again---and entered the novitiate of the Sisters
of Charity and Christian Instruction at Nevers, becoming Sister
Marie-Bernard. The \term{Litterae decretales} of her canonization
record the community's founding discipline about her past: the
day after her arrival, at the superior's order, she recounted all
the apparitions once, before the assembled sisters, ``candidly
and simply''---and thereafter never spoke of them again unless
obedience required it. Falling mortally ill in the novitiate, she
made her religious profession \term{in extremis}; recovered, she
completed it in regular form on 30 October 1867, and made
perpetual vows on 22 September 1878.\endnote{\work{Acta
Apostolicae Sedis} 26 (1934), pp.~76--77 (``ac deinceps de iisdem
apparitionibus, nisi obedientia adacta, haud amplius loquuta
est''). The Sanctuary's biography describes the same rule from
the community's side: she gave her account ``before the assembled
community on the day after she arrived; thereafter it was not to
be spoken of.''} Her assignment history was thin by
design---assistant infirmarian, sacristy work, and finally the
long assignment of being ill; on profession day, the bishop
reportedly declared that her work would be ``the work of
prayer.'' She understood her own position with the lucidity that
marks all her recorded sayings: ``I came here to hide myself'';
``My mission in Lourdes is finished''; and, writing to the Pope,
``My weapons are prayer and sacrifice.''\endnote{Sanctuary
biography page (received sayings, quoted as the Sanctuary
transmits them). Her sicknesses included the asthma that the
\term{Litterae decretales} say she bore ``with an even mind to
the last days of her mortal life'' (AAS 26 (1934), p.~76).}

She died at Nevers on 16 April 1879, aged thirty-five, having
received the sacraments. The \term{Litterae decretales} record
the immediate sequel with documentary soberness: the whole town
converged on the convent; the body remained supple and of natural
color for the four days of the viewing; she was buried in the
Saint-Joseph chapel within the convent enclosure, and the grave
quickly became a place of concourse and of reported
cures.\endnote{AAS 26 (1934), p.~77. This study repeats the
1933 acts' own description and asserts nothing further about the
preservation of the body or its later exhumations, which it did
not verify at source.}

\subsection{Beatification and canonization}

The cause proceeded at Rome through the Sacred Congregation of
Rites: Pius X signed the introduction of the cause on 13 August
1913; Pius XI declared her heroic virtue on 18 November 1923,
approved two beatification miracles by decree of 1 May 1925 (the
instantaneous cures of Henri Boisselet, of acute tubercular
disease, and of Sr Melania Maria Meyer of the Congregation of
Divine Providence, of a stomach ulcer), and beatified her in the
Vatican Basilica on 14 June 1925, in the holy year. For the
canonization two further cures were selected and judged: the
instantaneous cure of Alexis Lema\^itre, Archbishop of Carthage,
of long-standing severe amoebiasis, at Nevers during the
solemnities of the Beata's translation; and the cure of Sister
Marie de Saint-Fid\`ele (Amata Baug\'e) of the Good Shepherd
sisters \emph{at Lourdes}, of tubercular disease of the spine and
knee, during a novena to Blessed Bernadette in February 1928.
Pius XI issued the decree on the two miracles on 31 May 1933 and
the \term{tuto} decree on 2 July 1933, held the secret
consistory on 16 October, and canonized her on 8 December
1933---the feast of the Immaculate Conception, in the basilica
where the dogma had been defined seventy-nine years
before.\endnote{\work{Acta Apostolicae Sedis} 26 (1934),
pp.~78--80 (the \term{Litterae decretales}' procedural
narrative) and pp.~5--7 (the canonization rite of 8 December
1933). The cause chronology (1913 introduction under Pius X; 1923
heroic-virtue decree; 1 May 1925 miracles decree with the
Boisselet and Meyer cures) is at pp.~77--78; the OCR of the
volume is imperfect at the Boisselet diagnosis, which is
reported summarily here.}

\begin{operativetext}{Canonization formula of Pius XI, 8 December
1933 (Latin, as published in the Acta Apostolicae Sedis)}
Ad honorem Sanctae et individuae Trinitatis, ad exaltationem
Fidei catholicae et christianae Religionis augmentum, auctoritate
Domini Nostri Iesu Christi, Beatorum Apostolorum Petri et Pauli
ac Nostra; matura deliberatione praehabita et divina ope saepius
implorata, ac de Venerabilium Fratrum Nostrorum S.~R.~E.
Cardinalium, Patriarcharum, Archiepiscoporum et Episcoporum in
Urbe exsistentium consilio, Beatam \textsc{Mariam Bernardam
Soubirous Sanctam} esse decernimus et definimus, ac Sanctorum
catalogo adscribimus; statuentes ab Ecclesia universali eius
memoriam quolibet anno, die natali illius, nempe die XVI Aprilis,
inter Sanctas Virgines pia devotione recoli
debere.\endnote{\work{Acta Apostolicae Sedis} 26 (1934), p.~6
(the volume's OCR reads ``ad exarationem Fidei catholicae''; the
standard formula reads \term{exaltationem}, and the reading above
follows the formula with the OCR variant recorded here). In
editorial rendering: ``\ldots we decree and define that Blessed
Marie-Bernard Soubirous is a Saint, and we inscribe her in the
catalogue of the Saints, establishing that her memory is to be
devoutly kept by the universal Church every year on her heavenly
birthday, 16 April, among the Holy Virgins.''}
\end{operativetext}

\subsection{A saint of humility---what the canonization does and
does not establish}

The Sanctuary's biography closes with a sentence that this study
adopts as exactly right in both halves: ``The Church proclaimed
her a saint on 8th December, 1933, not for having been chosen for
the Apparitions, but for the way in which she responded to that
grace.''\endnote{Sanctuary biography page. The canonization
homily's exhortation identifies what is proposed for imitation:
``let us imitate her Christian lowliness of soul and her
humility; her faith; the ardent charity with which she burned''
and ``her zeal for penance'' (\work{Acta Apostolicae Sedis} 26
(1934), pp.~8--9, editorial translation of ``Eius nempe
christianam animi demissionem humilitatemque imitemur; eius
fidem, impensam eius, qua aestuabat, caritatem\ldots\ eius tamen
paenitentiae studium\ldots\ persequi contendamus'').} A
canonization is an act about a person's holiness and her
intercession, resting on heroic virtue and approved miracles; it
neither re-issues nor upgrades the 1862 apparition judgment, and
it does not make her every recollection inerrant. The Church's
current norms state the general principle expressly: even a
\term{Nihil obstat} granted in a canonization process ``does not
imply a declaration of authenticity regarding any supernatural
phenomena present in a person's life.''\endnote{DDF, \work{Norms
for Proceeding in the Discernment of Alleged Supernatural
Phenomena} (17 May 2024), n.~13, English text (References),
citing the canonization decree of St Gemma Galgani as its
example. The principle is stated by the norms as a general one;
its application to Bernadette's 1933 canonization is this
study's, and is consistent with the 1933 acts, which canonize her
virtue while narrating the apparitions as the received history of
her life.} That said, the 1933 acts do not bracket Lourdes: the
homily glories in it, and the canonization was staged on the
dogma's feast in the dogma's basilica (\S3.3). The Church of 1933
was not distancing itself from the apparitions; it was locating
their visionary's greatness where the Gospel locates greatness.
Pius XI's homily makes the point in its choice of texts: God
chose ``what is foolish in the world to shame the wise''---an
unlettered miller's daughter, ``rich in nothing but the candor of
her most sweet soul''---and the enduring Lourdes of the
sanctuaries and crowds is, in the homily's words, at once the
glory of the Immaculate Virgin and ``the singular monument of the
proven holiness of Marie-Bernard Soubirous.''\endnote{\work{Acta
Apostolicae Sedis} 26 (1934), pp.~7--8 (1~Cor 1:27 quoted;
``indoctam pistorum filiam, nulla re alia, nisi suavissimi animi
sui candore, praedivite''; ``Lapurdum, en suspicite\ldots\ Mariae
Bernardae Soubirous singulare comprobatae sanctitatis
monumentum'').}

Her liturgical memory was fixed by the canonization formula on 16
April; in France she is widely kept on 18 February. This study
records the universal date from the formula itself and the French
usage as common knowledge of the calendars, not re-verified in a
typical edition (Appendix~A).\endnote{Universal date: the formula
quoted above. The French date is a particular-calendar usage
reported at ceiling; no particular-law act was consulted.}
